% GOVERNANCE NOTE::LSPV2-WRITE-RECOVERY::2026-10-05T11:38:00-07:00
\documentclass[11pt]{article}
\usepackage[margin=.88in]{geometry}
\begin{document}
\section*{LSPV-2 consecutive write-failure recovery}
\textbf{STATUS: PROCEDURE REPAIR ONLY.} No scientific claim is changed. P versus NP remains open.

\paragraph{Evidence.}
LSPV-3 predecessor preflight at 10:34 PT found no verified default-branch LSPV-2 artifact attributable to the 10:22 PT scheduled run and issued a recovery instruction. The next LSPV-3 preflight at 11:37 PT again found no verified LSPV-2 artifact attributable to either the 10:22 or 11:22 PT runs. Repository search of the canonical LSPV-2 report directory returned older verified reports but no new 2026-10-05 10:22 or 11:22 report.

\paragraph{Classification.}
The exact failed scientific bodies and connector errors are not available to LSPV-3, so the scientific content must not be reconstructed. The repeated failure is classified as UNKNOWN at the write layer, with a possible procedure/state component. No repository permission failure is inferred.

\paragraph{Smallest repair.}
The LSPV-2 automation prompt was amended only at the write procedure: use the default branch by omitting an explicit branch; search the canonical path before retry; preserve an exact failed body if available; otherwise perform a fresh derivation under a new timestamp; after creation, immediately read the exact file back and require the report header and final verification sentinel; then reconcile only the LSPV-2 owned root region against the current root blob SHA. Placeholders and reconstruction from summaries remain forbidden.

\paragraph{Rollback.}
Remove the added write-procedure paragraph if verified LSPV-2 writes resume and repository governance later makes the repair redundant. The LSPV-2 schedule, enabled state, title, and scientific target were not changed.

% END_OF_VERIFIED_GOVERNANCE::LSPV2-WRITE-RECOVERY
\end{document}
